\section{Conclusion}\label{sec:conclusion}

Under drift, memory is never purely beneficial. It reduces finite-sample noise,
but it also ages. The resulting trade-off makes finite memory a
temporal-validity object with a finite horizon.

The horizon law
\[
  \mathbb E\,d\!\left(\widehat P_t^{(n)},P_t\right)
  \le C_K n^{-a} + C_S\zeta n^H,
\]
and its optimizer
\[
  n^*(a,H) \asymp (C_K/\zeta)^{1/(a+H)}
\]
summarize that balance. This law induces a useful-memory region around the
optimal scale.

The core theory closes in a tractable one-dimensional root-$n$ proof model,
matched with a structural Gaussian lower bound and a Gaussian location
benchmark. In fixed-$\varepsilon$ Sinkhorn geometry, the one-dimensional
inheritance results place the horizon law in an operational transport geometry.

The empirical signatures align with the theory: finite-memory U-curves, roughness-
dependent horizon scaling, a real-stream retention pattern, and a positive lag
between validity loss and detector alarm. Inferable horizons then add a second
distinction: a temporal-validity horizon can exist structurally and still be
operationally hard to estimate from lag geometry alone.

The open frontier consists of embedded Sinkhorn closure, non-uniform memory
design, and online control.
